The NMPC is developed using Acados, a software package that provides solvers for nonlinear optimal control.
It is based on the symbolic framework CasADi, and provides a Python interface at an higher level.
CasADI is a software framework for automatic differentiation, nonlinear optimization and 
optimal control. It is implemented in C++ and has also a Python interface.\\
The neural model has been developed using PyTorch, an open source deep learning library based on Python.\\
The linking between the MPC and the neural model has been possible using the L4Casadi framework, which allows
the integration of Pytorch-learned models within CasADI and Acados, and provides also an hardware acceleration
for learned components.\\
Acados was used for three main purposes:
\begin{itemize}
  \item Modeling
  \item Simulation through numerical integration
  \item Formulation and solving of optimization problems within MPC
\end{itemize}
The software of this project is based mainly on three components:
\begin{description}
    \item[Simulator] Matplotlib plot to show everything graphically in real-time. 
    \item[Model] Expressed formally according to Acados syntax, the ODE is expressed in implicit form.
    \item[Controller] Expressed as an Acados linear program, with a linear cost function and linear constraints; 
                    it returns a single control action based on the current state and the predicted state for
                    the whole horizon.
\end{description}

\subsection{Simulator}
The main goal of the simulator is to visualize the reference trajectory that evolves in time,
and to simulate the unicycle movement through the control inputs and integration of the model.
The complete parametrized implementation of a simulation can be found at 
\href{https://github.com/neverorfrog/simple_neural_mpc/blob/mlp-tests/src/simple_neural_mpc/simulation/trajectory_tracking.py}{trajectory\_tracking.py},
which takes in input the interested robot (in our case the unicycle, but it can be extended to other robots),
the used controller and the reference trajectory, and it generates the entire simulation loop. 
The simulation loop is summarized in pseudocode \ref{sim}, where the whole cycle is encapsulated in the function 
\begin{verbatim}
    run(self, N, animate, save)
\end{verbatim}
which takes the total number of steps $N$, a boolean to express if the simulation must be visualized and animated 
through a proper function \texttt{animate}, and another boolean if the animation video must be saved.
The animation is guaranteed by \texttt{FuncAnimation} of Matplotlib. \\
The trajectories the robot can track are defined inside 
\href{https://github.com/neverorfrog/simple_neural_mpc/blob/mlp-tests/src/simple_neural_mpc/utils/trajectory.py}{trajectory.py}.
A trajectory must contain a function \texttt{update(self, t)} which declares at time-step $t$
its position, velocity, acceleration, orientation and angular velocity.  
\begin{algorithm}[h]
    \caption{Simulation Loop}\label{sim}
    \begin{algorithmic}[1]
        \State $dt \gets 0.05$ s \Comment{Time step}
        \While{ number of steps is not exceeded }
        \State measure car state $x$
        \State action $a \gets command(x)$ \Comment{From MPC}
        \State move forward for $dt$ under action $a$, updating state $x$
        \State update plot and save state and action
        \EndWhile
    \end{algorithmic}\
\end{algorithm}

\subsection{Model}
An important aspect is how the models are actually discretized. What happens is that,
conceptually, the model is still expressed as a symbolic function through CasADI, as
if it was a continuous function. Acados helps in doing this, providing an high
level interface through the class \texttt{AcadosModel}. Instantiating an object of this class,
we have to deliver it state, inputs and derivative of the state as CasADI symbolic variables,
and also implicit or explicit ODE in the form of CasADI expression, according to the integrator
type wants to be used. All these information are necessary to the class to automatically
generate the CasADI function. The formulation of the real unicycle model according to 
Acados policy can be found in 
\href{https://github.com/neverorfrog/simple_neural_mpc/blob/mlp-tests/src/simple_neural_mpc/models/unicycle_kin/unicycle_kin_acados.py}{unicycle\_kin\_acados.py}.\\
In the same way, the neural model of the robot is expressed through the \texttt{AcadosModel}.
The only distinction lies in the expression of the explicit (or implicit) ODE, which recalls 
the neural approximator thanks to \textbf{L4Casadi}. In fact, it is possible to create 
an instance of the class \texttt{L4CasADi}, which takes in input the Pytorch tensor; 
at this point the forward pass of the neural network can be executed, calling the L4CasADi 
instance with the symbolic variables as inputs. In this way, a symbolic CasADi expression is returnes,
representing the explicit dynamics of the system. The detailed formulation of the neural unicycle
model can be found in 
\href{https://github.com/neverorfrog/simple_neural_mpc/blob/mlp-tests/src/simple_neural_mpc/models/unicycle_kin/unicycle_kin_acados_neural.py}{unicycle\_kin\_acados\_neural.py}  

\subsection{MPC}
\todo[inline]{da fare}